\paragraph{Conjugate-coordinate class: the \(\Theta\)-built heat kernel and the induced coefficient (conditional; bounded no-go).}
Keep the Poisson limit, the leaf \(\mathbb{CP}^{1,2}\), \(X^a=T^a+A^a\), the bilinear term with \(\mu^2=-3\) and the notation of the preceding paragraphs; \(D^a=\{T^a,\cdot\}\), \(J^{AB}=\{M^{AB},\cdot\}\), \(\Box=\eta_{ab}D^aD^b\). The regulator is the proper-time cutoff \(s\ge\LUV^{-2}\) of the operator \(\Box_X\) itself, applied to every sector; it is built from Poisson brackets and the matrix contraction only. Background gauge fixing \(D_aA^a\) gives the ghost operator \(-\Box\) (weight \(-2\)), the vector operator \(H_V=-\Box-2sJ_{ab}\Sigma^{ab}\) (the same \(J_{ab}A^b\) as in the quadratic form above, with \(\mu^2\) added), and, for \(n_{\mathrm{ex}}\) further matrices with trivial background, \(-\Box+\mu_i^2\) with \(\mu_i^2\) an additional input. Fermions are absent. In exact arithmetic in the polynomial Poisson algebra, for both signs \(s\): \(\Box=C_{\mathfrak{so}(4,1)}-C_{\mathfrak{so}(3,1)}\) with the generators \(J^{AB}\), \(A,B\le4\), and \(J^{ab}\), \(a,b\le3\); \([\Box,J^{ab}]=0\) while \([\Box,D^a]\ne0\); and on functions of \(y^A=M^{A5}\), \(A\le4\), one has \(D^a=s\,(y^4\partial^a-y^a\partial_4)\), a rotation field of the hyperboloid \(y\cdot y=-\ell^2\). Fibre-constant functions therefore form an invariant subspace of \(\Box\), on which \(\Box=\rho^{-1}\partial_\mu(\rho\,m^2\eta^{\mu\nu}\partial_\nu\,\cdot)\) with \(m=y^4\), \(\rho=1/m\); for \(G^{\mu\nu}=m^2\eta^{\mu\nu}\) the Seeley--DeWitt density of this sector is \(\sqrt{|G|}\,(R_G/6-E)=(5m^2-3\ell^2)/(4m^6)\), which is not of definite sign. Fibre harmonics of order \(\ge1\) are not an invariant subspace of the perturbed operators, so this sector does not by itself give the response to the level-one amplitudes.

\emph{Hypoelliptic structure.} At every sampled point of the orbit (random points, and the sequence \(y^4\to0\)) the fields \(D^a\) have rank four and \(D^a\), \(J^{ab}\) together span the six-dimensional tangent space; the horizontal rank drops on \(y^4=0\) (one singular value proportional to \(y^4\)), which is excluded. Hence \(\Box\) is a sub-Laplacian of step two with homogeneous dimension \(4+2\cdot2=8\), not an elliptic operator on a product. Under equiregularity the diagonal of its heat kernel is \(K(s)=c_0s^{-4}(1+O(s^{1/2}))\) with \(c_0\) given by the nilpotent approximation~\cite{BenArous1989}, here a group with four horizontal generators and a two-dimensional centre. After Fourier transform along the centre (variable \(\zeta\in\mathbb R^2\)) the nilpotent operator is a four-dimensional charged Laplacian in a constant field \(F_{ab}(\zeta)=\zeta_kS^{ab}{}_k\) of the two-plane strengths \(f_1,f_2\propto|\zeta|\), with Landau levels \(f_1(2n_1+1)+f_2(2n_2+1)\) and the diagonal \(\prod_j f_j/(4\pi\sinh sf_j)\) per unit horizontal volume and per \(\mathrm d^2\zeta/(2\pi)^2\). With \(\nu\) the Liouville volume of the frame \((D^0,\dots,D^3,Z_1,Z_2)\) of the nilpotent algebra, \(c_0=\nu^{-1}(64\pi^4)^{-1}\!\int_0^{2\pi}\!\mathrm d\varphi\,\phi_1\phi_2\!\int_0^\infty\!\mathrm du\,u^3/(\sinh u\phi_1\sinh u\phi_2)\), \(f_j=|\zeta|\phi_j\); for \(f_1=f_2\) this is \(3\zeta(3)/(64\pi^3\nu)\). With the Euclidean contraction \(\eta\to\delta\) of the four matrix indices (the Poisson algebra kept) the ratio \(f_1/f_2\) at the sampled points lies between \(0.006\) and \(0.84\), never at \(1\); for the Lorentzian contraction \(F\eta\) has one real and one imaginary eigenvalue pair, the electric-type case, for which no heat kernel exists.

\emph{Consequences.} (i) The statement that a fuzzy factor replicates the four-dimensional coefficient of Theorem~\ref{thm:Geff} over its modes, with threshold weights, is the limit \(sf_j\ll1\) of the diagonal above, \(\prod_j sf_j/\sinh sf_j=1-s^2(f_1^2+f_2^2)/6+\dots\), valid when \(\LUV^2\) exceeds the largest populated \(f_j\). In the Poisson limit \(|\zeta|\) is unbounded and the regime is the hypoelliptic one; the replication formula does not apply there, and the weighted mode count is replaced by \(\nu^{-1}\int\mathrm d^2\zeta\,(2\pi)^{-2}\prod_j sf_j/\sinh sf_j\). (ii) In the same nilpotent limit \(H_V\) acquires the Zeeman term \(\pm2f_1,\pm2f_2\), so its lowest eigenvalue at fixed \(\zeta\) is \(-|f_1-f_2|\) (checked by diagonalisation in a truncated Fock space): zero modes of unbounded number for \(f_1=f_2\), and eigenvalues unbounded below for \(f_1\ne f_2\), neither regulated by the proper-time cutoff. The vector sector therefore needs a fibre truncation (a finite fuzzy fibre) or fermionic partners, and in the Euclidean continuation used here it has no sign-definite heat trace. (iii) For \(sf_j\ll1\), the vector, ghost and \(n_{\mathrm{ex}}\) scalar sectors together give \(2+n_{\mathrm{ex}}+\tfrac{22-n_{\mathrm{ex}}}{6}s^2(f_1^2+f_2^2)\); the \(s^2\) term has the structure of the quartic term \(\{X^a,X^b\}^2\) and renormalises its coefficient logarithmically, not the Einstein--Hilbert coefficient. (iv) An Einstein--Hilbert coefficient would be the weight-two term \(\propto s^{-3}\) of the diagonal and its second variation under \(A\); no such computation is made here, and neither the sign nor the size of \(w\), hence of \(\varrho=\Lambda/w\), follows.

This is a statement for the Poisson limit, the leading (nilpotent) symbol, equiregular points, the Euclidean contraction \(\eta\to\delta\) with the Lie--Poisson algebra kept, background gauge and the bosonic content listed. It does not treat Moyal terms, fermions, the subleading symbol, the dependence on \(A\), or the identification of \(G\) with the emergent metric of Chapter~\ref{ch:vierbein-emergence}, and it does not bear on Assumption~\ref{ass:EE}. This changes no ledger label; Problem~F remains open.

